\documentclass[]{deedy-resume-openfont}
\usepackage{fancyhdr}

\pagestyle{fancy}
\fancyhf{}

\begin{document}

\namesection{尹}{达恒}{
  \urlstyle{same}\href{mailto:dya64@sfu.ca}{dya64@sfu.ca}
  | Simon Fraser University, Burnaby, Canada
}

\begin{minipage}[t]{0.24\textwidth}

\section{教育经历}
\subsection{Simon Fraser Univ.}
\descript{计算科学 博士在读}
\descript{导师：刘江川}
\location{2023.08 -- 至今}
\sectionsep

\subsection{东南大学}
\descript{计算机科学与技术 硕士学位}
\descript{导师：东方}
\location{2020.08 -- 2023.07}
\sectionsep

\subsection{江南大学}
\descript{物联网工程 学士学位}
\descript{导师：郭亚}
\location{2016.08 -- 2020.07}
\descript{GPA: 3.63, Rank 9/141}
\sectionsep

\section{研究方向}
\textbullet{} 体积视频传输 \\
\textbullet{} 动态三维重建 \\
\textbullet{} 具身智能系统 \\
\sectionsep

\section{链接}
Homepage: \href{https://www.yindaheng98.top}{\bf yindaheng98.top} \\
Github: \href{https://github.com/yindaheng98}{\bf @yindaheng98}
\sectionsep

\end{minipage}
\hfill
\begin{minipage}[t]{0.74\textwidth}

\section{研究概述}
\sectionsep
\begin{tightemize}
\item 我致力于构建连接真实物理环境与Physical AI的动态三维实景系统：通过实时感知与\\三维重建，将不断变化的真实环境转化为可远程访问的动态三维实景上下文，并在网络\\波动、带宽受限、边缘算力不足的真实部署约束下，仍能为具身智能等Physical AI\\系统提供低时延的物理世界感知基础。
\item 近期工作围绕动态三维实景的重建、压缩与弱网传输展开，基于动态3D Gaussian \\Splatting、通过点追踪运动先验、自适应LoD表示与客户端侧视觉恢复等技术，在\\网络波动和边缘算力受限条件下保持实景系统的稳定性、清晰度与交互质量。
\end{tightemize}
\sectionsep

\section{代表性论文}
\runsubsection{CAGS: Color-Adaptive Volumetric Video Streaming with Dynamic 3D Gaussian Splatting}
\descript{SIGGRAPH, 2026}
\begin{tightemize}
  \item \textbf{Daheng Yin}, Yili Jin, Jianxin Shi, Isaac Ding, Miao Zhang, Fangxin Wang, Zhaowu Huang, Cong Zhang, Jiangchuan Liu, Fang Dong
  \item 核心思想：提出面向动态3DGS体视频流传输的色彩自适应机制，以向量量化构建多层\\级LoD传输表示，并利用低分辨率参考图在客户端修正量化引起的颜色失真，从而在\\带宽波动下实现更高质量、更低开销的自适应传输。
\end{tightemize}
\sectionsep

\runsubsection{TrackerSplat: Exploiting Point Tracking for Fast and Robust Dynamic 3D Gaussians Reconstruction}
\descript{SIGGRAPH Asia, 2025}
\begin{tightemize}
  \item \textbf{Daheng Yin}, Isaac Ding, Yili Jin, Jianxin Shi, Jiangchuan Liu
  \item 核心思想：利用点跟踪先验引导高斯椭球运动，提升动态场景重建的稳健性与效率。
\end{tightemize}
\sectionsep

\runsubsection{FSVFG: Towards Immersive Full-Scene Volumetric Video Streaming with Adaptive Feature Grid}
\descript{ACM Multimedia, 2024}
\begin{tightemize}
  \item \textbf{Daheng Yin}, Jianxin Shi, Miao Zhang, Zhaowu Huang, Jiangchuan Liu, Fang Dong
  \item 核心思想：以特征网格表示全景体积视频，并按带宽自适应传输关键特征与帧间\\残差，实现高效体积视频流传输。
\end{tightemize}
\sectionsep

\section{科研与工程经历}

\runsubsection{Fast and Robust Dynamic 3D Gaussian Reconstruction}
\descript{2024 -- 至今}
\begin{tightemize}
  \item 设计基于点跟踪的高斯运动补偿方法，解决动态场景帧间位移导致的重建退化问题。
  \item 搭建多视角多GPU并行重建流程，在保持时序一致性的同时提升重建吞吐。
  \item 面向下游任务增强运动线索可靠性，支持具身场景理解与仿真相关应用。
\end{tightemize}
\sectionsep

\runsubsection{Adaptive Full-Scene Volumetric Video Streaming}
\descript{2023 -- 至今}
\begin{tightemize}
  \item 构建基于特征网格与高斯表示的全场景动态体视频自适应传输系统。
  \item 设计LoD、残差调度与ABR机制，在真实网络波动下优化画质与时延平衡。
\end{tightemize}
\sectionsep

\runsubsection{Collaborative Live Video Super-Resolution}
\descript{2021 -- 2023}
\begin{tightemize}
  \item 开展分布式边缘推理调度与低时延多媒体传输研究。
  \item 为后续体视频重建与传输方向提供了系统基础与工程经验。
\end{tightemize}
\sectionsep

\section{竞赛与荣誉}
\begin{tabular}{lll}
2025 & SIGGRAPH Asia Volumetric Video Workshop & 视频压缩方向 3rd \\
2022 & TensorRT Hackathon (NVIDIA \& TIANCHI) & 优胜奖 \\
2021 & TensorRT Hackathon (NVIDIA \& TIANCHI) & 排名 4/48 \\
2020 & 江南大学本科优秀毕业设计 & \\
2018 & 江南大学2017-2018学年一等学业奖学金 &  \\
2018 & 2018年全国大学生数学建模竞赛 & 国家级二等奖 \\
2017 & 第九届全国大学生数学竞赛 & 江苏赛区二等奖 \\
2017 & 2016-2017年度国家奖学金 & \\
2017 & 江苏省普通高等学校第十四届数学竞赛 & 本科一级组一等奖 \\
2017 & 入选江南大学至善学院（荣誉学院）& \\
\end{tabular}
\sectionsep

\end{minipage}
\end{document}
